% !TeX root = ../../main.tex
\subsection{A filtration by odd Chern characters}
\label{subsec:s05:01}

We use the boundary exponential in
\cite[Lemma~3.10, (3.18)--(3.21)]{ABCDGLW23} for a general symmetric
node series. Its scalar specialization was proved in
Proposition~\ref{prp:s02:boundary-exponential}. We state the labeled
specialization and give its algebraic proof.

Fix an integer $r\geq1$ and odd integers $a_1,\ldots,a_r\geq3$,
and put $w=\sum_{i=1}^r a_i$. For an odd integer $a\geq3$ and formal variables $x,y$, define
the symmetric polynomial
\begin{equation}\label{eq:s05:node-polynomial}
 \mathcal P_a(x,y)=\sum_{c=0}^{a-1}(-1)^c x^c y^{a-1-c}.
\end{equation}
Let $\varepsilon_1,\ldots,\varepsilon_r$ be commuting formal
variables with $\varepsilon_i^2=0$. Coefficient extraction refers to
this square-free algebra. Set
\begin{equation}\label{eq:s05:source-kernel}
\begin{split}
 \mathcal E(x,y)&=\alpha+
       \sum_{i=1}^r\varepsilon_i\mathcal P_{a_i}(x,y),\\
 K_{\mathcal E}(x,y)&=
       \frac{1-e^{-(x+y)\mathcal E(x,y)}}{x+y}.
\end{split}
\end{equation}
The quotient is a formal series: its numerator is divisible by $x+y$.
All Chow series in the next identity are truncated by codimension.

\begin{lmm}[Labeled sources]\label{lmm:s05:labelled-sources}
For $g\geq2$,
\begin{align}
 &\int_{\MM_g}e^{12\alpha\lambda_1}
              \prod_{i=1}^r\mathrm{ch}_{a_i}\notag\\
 &=\left(\prod_{i=1}^r\frac{B_{a_i+1}}{(a_i+1)!}\right)
 [\varepsilon_1\cdots\varepsilon_r]
 \sum_{\Gamma\in\mathcal G_{g,0}}
 \frac1{|\operatorname{Aut}\Gamma|}\notag\\
 &\qquad\times\int_{\MM_\Gamma}
 \prod_{v\in V(\Gamma)}
  \exp\left(\alpha\kappa_1^{(v)}+
          \sum_{i=1}^r\varepsilon_i\kappa_{a_i}^{(v)}\right)
 \prod_{\substack{e\in E(\Gamma)\\e=\{f,f'\}}}
       K_{\mathcal E}(\psi_f,\psi_{f'}).
 \label{eq:s05:source-graphs}
\end{align}
\end{lmm}
\begin{prf}
On the unmarked space, \eqref{eq:s02:grr} writes each character as
$B_{a+1}/(a+1)!$ times
$\kappa_a+\tfrac12\xi_*\mathcal P_a(\psi_+,\psi_-)$.
Also $12\alpha\lambda_1=\alpha\kappa_1+\alpha\delta$.
The coefficient of the product of all the sources in the
exponential of these classes is their product, with no $r!$.

The normal line of a boundary branch has first Chern class
$-(x+y)$. Repeating that branch $k\geq1$ times in the exponential
therefore gives
\[
 \sum_{k\geq1}\frac{(-(x+y))^{k-1}\mathcal E(x,y)^k}{k!}
 =K_{\mathcal E}(x,y).
\]
Distinct nodes give the product of these factors, exactly as in the
proof of Proposition~\ref{prp:s02:boundary-exponential}. The
logarithmic dualizing line restricts to the logarithmic dualizing
line on each component, with the branches included among its
markings. Integrating its powers on the normalized universal curve
thus gives $\xi_\Gamma^*\kappa_a=\sum_v\kappa_a^{(v)}$.
The projection formula supplies the displayed vertex exponentials.

Finally, passing from labeled vertex and flag products to graph
isomorphism classes gives $1/|\operatorname{Aut}\Gamma|$.
The source labels remain distinct. Equivalently, grouping their
assignments into orbits gives the stabilizer of each decorated
graph; no permutation factor for the source labels is introduced.
All steps are self-intersection and pushforward operations in Chow,
so integration proves \eqref{eq:s05:source-graphs}.
\end{prf}

\begin{rmk}[Boundary normalization]
The cited lemma is stated in cohomology. Its normal-line proof,
spelled out above in the required specialization, is algebraic.
The source's boundary pushforward is the sum over one-edge graphs
with their automorphism factors. This equals our
$\tfrac12\xi_*$ on symmetric branch classes, already used in
\eqref{eq:s02:ordered-boundary}; no additional factor $1/2$ belongs
in \eqref{eq:s05:source-kernel}.
\end{rmk}

A source is assigned either to an edge or to a vertex. If a nonempty set
$J\subseteq[r]$ is assigned to one edge, its coefficient is
\begin{equation}\label{eq:s05:tagged-edge}
 [\textstyle\prod_{i\in J}\varepsilon_i]K_{\mathcal E}(x,y)
 =(-1)^{|J|-1}(x+y)^{|J|-1}e^{-\alpha(x+y)}
       \prod_{i\in J}\mathcal P_{a_i}(x,y).
\end{equation}
The polynomial multiplying the exponential has degree
$\sum_{i\in J}a_i-1$. An edge with no source retains
$K_\alpha(x,y)=\int_0^\alpha e^{-t(x+y)}\dd t$.

If a set $I\subseteq[r]$ is assigned to a vertex of type $(G,n)$,
use the partition inversion \eqref{eq:s02:partition-inversion}.
For each partition $\mathcal A$ of $I$, order its blocks as
$C_1,\ldots,C_m$, where $m=|\mathcal A|$; the expression is
independent of this order. For any polynomial, or
codimension-truncated series, $Q$ in the old cotangent classes,
\begin{align}
 &\int_{\MM_{G,n}}e^{\alpha\kappa_1}Q
                   \prod_{i\in I}\kappa_{a_i}\notag\\
 &=\sum_{\mathcal A\text{ a partition of }I}(-1)^{|I|-m}
   \int_{\MM_{G,n+m}}
    e^{\alpha\kappa_1-\alpha\sum_{k=1}^m\psi_{n+k}}
    Q
    \prod_{k=1}^m\psi_{n+k}^{\,1+\sum_{i\in C_k}a_i}.
 \label{eq:s05:free-markings}
\end{align}
For $I$ empty this means the original integral. Every new cotangent
power is at least four. Lemma~\ref{lmm:s02:rational-tail-degree}
therefore replaces the pullbacks of $Q$ by the intrinsic old classes
and replaces the pulled-back background by the exponential displayed
in \eqref{eq:s05:free-markings}. In particular the new factors
$e^{-\alpha\psi_{n+k}}$ survive. The inverse has the sign shown and
has no block factorials.

Each free new marking can return only one power of $g$ against the
loss $g^{-a}$ from degree conversion for a character of degree $a$.
We now justify this count for the complete graph sum, before using
it to discard any terms.

\begin{lmm}[Character filtration]\label{lmm:s05:character-filtration}
For the fixed indices above, as $g\to\infty$,
\begin{equation}\label{eq:s05:character-bound}
 \left\langle\prod_{i=1}^r\mathrm{ch}_{a_i}\right\rangle_g
       =O_{a_1,\ldots,a_r}(g^{r-w}).
\end{equation}
\end{lmm}
\begin{prf}
\emph{Uniform bounds on the original graphs.}
Estimate each term of \eqref{eq:s05:source-graphs} on its original
unmarked graph.
Let $s_v$ be the number of forgetting blocks at vertex $v$ and
$m=\sum_v s_v\leq r$. All cotangent insertions in
\eqref{eq:s05:tagged-edge} and \eqref{eq:s05:free-markings} are
fixed derivatives of vertex functions. Their total derivative order
is bounded in terms of the fixed $a_i$. Applying the multivariable
Cauchy formula on fixed discs to
Lemma~\ref{lmm:s03:vertex-bound} bounds each such derivative by the
same vertex majorant, with changed constants independent of the
stable type. Edge parameters lie in $[0,\alpha]$ or equal $\alpha$;
the discs lie in one fixed compact set. At a vertex with $s_v$ new
markings the factorial part changes by at most
\[
 \frac{\Gamma(k_v+s_v)}{\sqrt{k_v+s_v}}
 \leq(k_v+r)^{s_v}\frac{\Gamma(k_v)}{\sqrt{k_v}},
 \qquad k_v=2g_v-2+n_v\geq1.
\]
The other new factors are bounded constants depending on the indices.

In this counting paragraph put $H=2g-2$. Summing the assignments of
$m$ labeled blocks to vertices costs at most
\[
 \left(\sum_v(k_v+r)\right)^m
 \leq((r+1)H)^m,
\]
since $\sum_v k_v=H$ and $|V(\Gamma)|\leq H$.
There are only finitely many partitions of the fixed source labels.
Assigning the remaining labels to at most $r$ specified edges costs
at most $C_r(1+|E(\Gamma)|)^r$. This is absorbed into
$T^{|E(\Gamma)|}$ for a fixed $T>1$, rather than into another power
of $H$. Thus, even when absolute values are taken before the finite
source and partition sums, the complete graph term is bounded by
\begin{equation}\label{eq:s05:source-majorant}
 \frac{C(1+H)^rT^{|E(\Gamma)|}}{|\operatorname{Aut}\Gamma|}
 \prod_{v\in V(\Gamma)}
 \frac{C_1B^{n_v}(4\pi^2)^{k_v-1}\Gamma(k_v)}{\sqrt{k_v}}.
\end{equation}
All constants are independent of $g$ and $\Gamma$.

For every fixed integer $M\geq1$,
Corollary~\ref{crl:s03:marked-extension}\textup{(ii)}, applied to
these original graphs with no external legs and polynomial loss $r$,
now bounds the absolute contribution of graphs with no large vertex
or with $q\geq M$ by
\begin{equation}\label{eq:s05:source-tail}
 O_{a_1,\ldots,a_r,M}(F_g g^{r-M}).
\end{equation}

\emph{Fixed cores and all legal loop counts.}
Here $q$ and the core refer to the original unmarked graph, before
adding forgetting markings. Fix one core, a source assignment and
its partitions, and let $s$ be the number of their new markings at
the large vertex, so $0\leq s\leq r$. Its complexity after these
markings is $2g-2-q+s$, independently of the restored loop count
$p$. The all-type bound of Lemma~\ref{lmm:s03:vertex-bound} gives
the factorial ratio
$\Gamma(2g-2-q+s)/\Gamma(2g-2)=O(g^{s-q})$.
The small vertices have fixed types, and choices of specified
restored loops add only a fixed polynomial in $p$. Therefore, for
constants $C>0$ and an integer $A\geq0$ depending on these fixed
data, the absolute contribution with any legal $p$ is at most
\begin{equation}\label{eq:s05:decorated-core-bound}
 F_g g^{s-q}\frac{C^{p+1}(1+p)^A}{p!}.
\end{equation}
The all-type bound applies even when the large vertex's genus is
small. Summing over legal $p>\log g$ gives a remainder smaller
than any prescribed inverse power. Only after that truncation may
one use HMT's expansions on $p\leq\log g$, as in
Proposition~\ref{prp:s03:core-expansion}. In particular the bound
summed over all legal $p$ is $O(F_g g^{s-q})$.

Take $M=1$ in \eqref{eq:s05:source-tail}. The remaining finitely
many core assignments have $q=0$ and $s\leq r$, so their total is
$O(F_g g^r)$. Together with the tail this bounds the entire
integral in \eqref{eq:s05:source-graphs} by $O(F_g g^r)$.
Finally, for $3g-3\geq w$, degree conversion gives
\begin{equation}\label{eq:s05:character-normalization}
 \left\langle\prod_{i=1}^r\mathrm{ch}_{a_i}\right\rangle_g
 =(12\alpha)^w\frac{(3g-3-w)!}{(3g-3)!}
 \frac{\displaystyle\int_{\MM_g}e^{12\alpha\lambda_1}
                    \prod_{i=1}^r\mathrm{ch}_{a_i}}
      {Z_g(1)}.
\end{equation}
The factorial prefactor is $O(g^{-w})$.
Theorem~\ref{thm:s03:graph-expansion} gives
$Z_g(1)\sim e^{S(1)}V_{g,0}$.
Since $V_{g,0}\asymp F_g$, this proves
\eqref{eq:s05:character-bound}.
\end{prf}

The Chern exponential now reduces the theorem to one character.

\begin{crl}[Algebraic reduction]\label{crl:s05:algebraic-reduction}
For the fixed product in Theorem~\ref{thm:fixed-insertions},
\begin{equation}\label{eq:s05:algebraic-reduction}
 F(j)\langle P(\lambda)\rangle_g
 =1+2G(j)\langle\mathrm{ch}_3\rangle_g+O(g^{-4}).
\end{equation}
\end{crl}
\begin{prf}
For a homogeneous Hodge polynomial $Q$ and an integer $c\geq0$
with $\deg Q+c\leq3g-3$, the definition of the normalized bracket
gives the exact identity
$\langle\lambda_1^c Q\rangle_g=\langle Q\rangle_g$.
In \eqref{eq:s02:chern-exponential}, the term with no character of
degree at least three contributes $\lambda_1^{j_k}/j_k!$ to
$\lambda_{j_k}$. The term with one $\mathrm{ch}_3$ contributes
$2\mathrm{ch}_3\lambda_1^{j_k-3}/(j_k-3)!$ when $j_k\geq3$.
Multiplying the fixed factors and then $F(j)$ gives the coefficient
$2\sum_{k=1}^{\ell}j_k(j_k-1)(j_k-2)=2G(j)$.

Every other nonconstant term has either one odd character of degree
at least five, or at least two characters of degrees at least three.
Lemma~\ref{lmm:s05:character-filtration} makes each such normalized
term $O(g^{-4})$. There are only finitely many terms, since the
$j_k$ are fixed. Indices below three simply supply no
$\mathrm{ch}_3$ term, completing the proof.
\end{prf}
